\documentclass[10pt,a4paper]{amsart}
\usepackage{amsmath,amssymb,amsthm,graphicx}
\usepackage{verbatim}
\newenvironment{reference}{\comment}{\endcomment}
%\newenvironment{reference}{}{}
\newenvironment{slogan}{\comment}{\endcomment}
\newenvironment{history}{\comment}{\endcomment}

% For commutative diagrams we use Xy-pic
\usepackage[all]{xy}

% We use 2cell for 2-commutative diagrams.
\xyoption{2cell}
% Theorem environments.
%
\theoremstyle{plain}
\newtheorem{theorem}[subsection]{Theorem}
\newtheorem{proposition}[subsection]{Proposition}
\newtheorem{lemma}[subsection]{Lemma}

\theoremstyle{definition}
\newtheorem{definition}[subsection]{Definition}
\newtheorem{example}[subsection]{Example}
\newtheorem{exercise}[subsection]{Exercise}
\newtheorem{situation}[subsection]{Situation}

\theoremstyle{remark}
\newtheorem{remark}[subsection]{Remark}
\newtheorem{remarks}[subsection]{Remarks}

\numberwithin{equation}{subsection}

% Macros
%
\def\lim{\mathop{\mathrm{lim}}\nolimits}
\def\colim{\mathop{\mathrm{colim}}\nolimits}
\def\Spec{\mathop{\mathrm{Spec}}}
\def\Hom{\mathop{\mathrm{Hom}}\nolimits}
\def\Ext{\mathop{\mathrm{Ext}}\nolimits}
\def\SheafHom{\mathop{\mathcal{H}\!\mathit{om}}\nolimits}
\def\SheafExt{\mathop{\mathcal{E}\!\mathit{xt}}\nolimits}
\def\Sch{\mathit{Sch}}
\def\Mor{\mathop{\mathrm{Mor}}\nolimits}
\def\Ob{\mathop{\mathrm{Ob}}\nolimits}
\def\Sh{\mathop{\mathit{Sh}}\nolimits}
\def\NL{\mathop{N\!L}\nolimits}
\def\CH{\mathop{\mathrm{CH}}\nolimits}
\def\proetale{{pro\text{-}\acute{e}tale}}
\def\etale{{\acute{e}tale}}
\def\QCoh{\mathit{QCoh}}
\def\Ker{\mathop{\mathrm{Ker}}}
\def\Im{\mathop{\mathrm{Im}}}
\def\Coker{\mathop{\mathrm{Coker}}}
\def\Coim{\mathop{\mathrm{Coim}}}

% Boxtimes
%
\DeclareMathSymbol{\boxtimes}{\mathbin}{AMSa}{"02}

%
% Macros for moduli stacks/spaces
%
\def\QCohstack{\mathcal{QC}\!\mathit{oh}}
\def\Cohstack{\mathcal{C}\!\mathit{oh}}
\def\Spacesstack{\mathcal{S}\!\mathit{paces}}
\def\Quotfunctor{\mathrm{Quot}}
\def\Hilbfunctor{\mathrm{Hilb}}
\def\Curvesstack{\mathcal{C}\!\mathit{urves}}
\def\Polarizedstack{\mathcal{P}\!\mathit{olarized}}
\def\Complexesstack{\mathcal{C}\!\mathit{omplexes}}
% \Pic is the operator that assigns to X its picard group, usage \Pic(X)
% \Picardstack_{X/B} denotes the Picard stack of X over B
% \Picardfunctor_{X/B} denotes the Picard functor of X over B
\def\Pic{\mathop{\mathrm{Pic}}\nolimits}
\def\Picardstack{\mathcal{P}\!\mathit{ic}}
\def\Picardfunctor{\mathrm{Pic}}
\def\Deformationcategory{\mathcal{D}\!\mathit{ef}}
\usepackage{cleveref}
\crefname{originalreference}{original reference}{original references}
\Crefname{originalreference}{Original reference}{Original references}
\makeatletter
\newlabel{situation-setup-base-change}{{0FVG}{source}{}{sourcepointer-1}{}}
\newlabel{situation-setup-base-change@cref}{{[originalreference][0][]0FVG}{[1][0][]source}{}{}{}}
\makeatother
\crefname{originallemma}{original lemma}{original lemmas}
\Crefname{originallemma}{Original lemma}{Original lemmas}
\makeatletter
\newlabel{morphisms-lemma-base-change-flat}{{01U9}{source}{}{sourcepointer-2}{}}
\newlabel{morphisms-lemma-base-change-flat@cref}{{[originallemma][0][]01U9}{[1][0][]source}{}{}{}}
\makeatother
\crefname{originallemma}{original lemma}{original lemmas}
\Crefname{originallemma}{Original lemma}{Original lemmas}
\makeatletter
\newlabel{morphisms-lemma-generalizations-lift-flat}{{03HV}{source}{}{sourcepointer-3}{}}
\newlabel{morphisms-lemma-generalizations-lift-flat@cref}{{[originallemma][0][]03HV}{[1][0][]source}{}{}{}}
\makeatother
\crefname{originallemma}{original lemma}{original lemmas}
\Crefname{originallemma}{Original lemma}{Original lemmas}
\makeatletter
\newlabel{schemes-lemma-fibre-topological}{{01K1}{source}{}{sourcepointer-4}{}}
\newlabel{schemes-lemma-fibre-topological@cref}{{[originallemma][0][]01K1}{[1][0][]source}{}{}{}}
\makeatother
\crefname{originallemma}{original lemma}{original lemmas}
\Crefname{originallemma}{Original lemma}{Original lemmas}
\makeatletter
\newlabel{morphisms-lemma-dimension-fibre-after-base-change}{{02FY}{source}{}{sourcepointer-5}{}}
\newlabel{morphisms-lemma-dimension-fibre-after-base-change@cref}{{[originallemma][0][]02FY}{[1][0][]source}{}{}{}}
\makeatother
\crefname{originallemma}{original lemma}{original lemmas}
\Crefname{originallemma}{Original lemma}{Original lemmas}
\makeatletter
\newlabel{morphisms-lemma-dimension-fibre-at-a-point}{{02FX}{source}{}{sourcepointer-6}{}}
\newlabel{morphisms-lemma-dimension-fibre-at-a-point@cref}{{[originallemma][0][]02FX}{[1][0][]source}{}{}{}}
\makeatother
\crefname{originaldefinition}{original definition}{original definitions}
\Crefname{originaldefinition}{Original definition}{Original definitions}
\makeatletter
\newlabel{definition-delta-dimension}{{02QP}{source}{}{sourcepointer-7}{}}
\newlabel{definition-delta-dimension@cref}{{[originaldefinition][0][]02QP}{[1][0][]source}{}{}{}}
\makeatother
\title{Base change shifts the delta-dimension of every irreducible component by the constant c}
\author{The Stacks Project Authors}
\date{stacks-project:chow.tex:local-snapshot-20261012; GNU FDL 1.2 or later}
\begin{document}
\maketitle
\noindent\textit{Editorial note.} All delivered mathematical text is the author's own, copied from the pinned Stacks Project chapter file chow.tex; no statement or proof step was rewritten, reordered or removed.\par\medskip
\noindent\textit{Editorial note.} Preamble provenance: the pinned chapter file opens with an input of the shared preamble and carries no preamble text of its own, so the project's single Stacks-wide preamble staged at sources/stacks-chow/preamble.tex is excerpted verbatim as two spans, lines 14-18 (Xy-pic setup) and lines 154-217 (the native theorem declarations and mathematical macros). Nothing was added, removed or redefined; the wrapper class and the amsmath/amssymb/amsthm/graphicx base packages are supplied by the assembler.\par\medskip
\noindent\textit{Editorial note.} The counted statement and proof contain no cross reference, no citation and no included file, so the delivered body is self contained.\par\medskip
\noindent\textit{Editorial note.} Canonical location of the counted proposition: Stacks tag 0FVH (sources/stacks-chow/pinned-original-source.tex line 15139).\par\medskip
\noindent\textit{Allowed standard prerequisites (editorial).}\begin{itemize}
\item Flatness of base change and lifting of generalizations along flat morphisms. Morphisms, Lemmas labelled morphisms-lemma-base-change-flat and morphisms-lemma-generalizations-lift-flat, cited at proof lines 15146-15153. Shows X' -> X is flat and that generic points of components of X' lie over the generic point of X.
\item Fibres and their topological description. Schemes, Lemma labelled schemes-lemma-fibre-topological, cited at proof lines 15157-15160. Identifies S'\_s with g\textasciicircum{}\{-1\}(s) so that x' can be seen as a point of S'\_s x\_s x.
\item Equality of fibre dimensions and of transcendence degrees after base change. Morphisms, Lemmas labelled morphisms-lemma-dimension-fibre-after-base-change and morphisms-lemma-dimension-fibre-at-a-point, cited at proof lines 15169-15175. Gives dim\_x(X\_s) = dim\_\{x'\}(X\_\{s'\}) and the equality of the two transcendence degrees.
\item The delta-dimension via the dimension function and the transcendence degree. Same chapter, Definition labelled definition-delta-dimension, cited at proof lines 15177-15179. Converts the displayed equality into the statement about dim\_\{delta'\}(Z').
\end{itemize}\medskip
\ifdefined\hypertarget\hypertarget{sourcepointer-1}{}\fi
\noindent\textit{Source reference.} Original reference 0FVG: sources/stacks-chow/pinned-original-source.tex, source lines 15123--15123. Same-chapter pointer into the author environment situation at source line 15123; invoked at proof line None of the counted proof. It is retained only as a navigation pointer; every load-bearing step is inside the delivered span.\par\medskip
\ifdefined\hypertarget\hypertarget{sourcepointer-2}{}\fi
\noindent\textit{Source reference.} Original lemma 01U9: sources/stacks-chow/tags.txt, source lines 2402--2402. Cross-chapter standard prerequisite located in the project tag list staged in the local snapshot; invoked at proof line 15148 of the counted proof. It is retained only as a navigation pointer; every load-bearing step is inside the delivered span.\par\medskip
\ifdefined\hypertarget\hypertarget{sourcepointer-3}{}\fi
\noindent\textit{Source reference.} Original lemma 03HV: sources/stacks-chow/tags.txt, source lines 4581--4581. Cross-chapter standard prerequisite located in the project tag list staged in the local snapshot; invoked at proof line 15152 of the counted proof. It is retained only as a navigation pointer; every load-bearing step is inside the delivered span.\par\medskip
\ifdefined\hypertarget\hypertarget{sourcepointer-4}{}\fi
\noindent\textit{Source reference.} Original lemma 01K1: sources/stacks-chow/tags.txt, source lines 2034--2034. Cross-chapter standard prerequisite located in the project tag list staged in the local snapshot; invoked at proof line 15156 of the counted proof. It is retained only as a navigation pointer; every load-bearing step is inside the delivered span.\par\medskip
\ifdefined\hypertarget\hypertarget{sourcepointer-5}{}\fi
\noindent\textit{Source reference.} Original lemma 02FY: sources/stacks-chow/tags.txt, source lines 3209--3209. Cross-chapter standard prerequisite located in the project tag list staged in the local snapshot; invoked at proof line 15166 of the counted proof. It is retained only as a navigation pointer; every load-bearing step is inside the delivered span.\par\medskip
\ifdefined\hypertarget\hypertarget{sourcepointer-6}{}\fi
\noindent\textit{Source reference.} Original lemma 02FX: sources/stacks-chow/tags.txt, source lines 3208--3208. Cross-chapter standard prerequisite located in the project tag list staged in the local snapshot; invoked at proof line 15172 of the counted proof. It is retained only as a navigation pointer; every load-bearing step is inside the delivered span.\par\medskip
\ifdefined\hypertarget\hypertarget{sourcepointer-7}{}\fi
\noindent\textit{Source reference.} Original definition 02QP: sources/stacks-chow/pinned-original-source.tex, source lines 1504--1504. Same-chapter pointer into the author environment definition at source line 1504; invoked at proof line 15178 of the counted proof. It is retained only as a navigation pointer; every load-bearing step is inside the delivered span.\par\medskip
\begin{lemma}
\label{lemma-dimension-base-change}
In Situation \ref{situation-setup-base-change} let $X \to S$ be locally
of finite type. Denote $X' \to S'$ the base change by $S' \to S$.
If $X$ is integral with $\dim_\delta(X) = k$, then
every irreducible component $Z'$ of $X'$ has $\dim_{\delta'}(Z') = k + c$,
\end{lemma}

\begin{proof}
The projection $X' \to X$ is flat as a base change of the flat morphism
$S' \to S$ (Morphisms, Lemma \ref{morphisms-lemma-base-change-flat}).
Hence every generic point $x'$ of an irreducible
component of $X'$ maps to the generic point $x \in X$ (because generalizations
lift along $X' \to X$ by
Morphisms, Lemma \ref{morphisms-lemma-generalizations-lift-flat}).
Let $s \in S$ be the image of $x$.
Recall that the scheme $S'_s = S' \times_S s$
has the same underlying topological space as $g^{-1}(\{s\})$
(Schemes, Lemma \ref{schemes-lemma-fibre-topological}).
We may view $x'$ as a point of the scheme $S'_s \times_s x$ which
comes equipped with a monomorphism $S'_s \times_s x \to S' \times_S X$.
Of course, $x'$ is a generic point of an irreducible component
of $S'_s \times_s x$ as well.
Using the flatness of $\Spec(\kappa(x)) \to \Spec(\kappa(s)) = s$
and arguing as above, we see that $x'$ maps to a generic point $s'$
of an irreducible component of $g^{-1}(\{s\})$. Hence
$\delta'(s') = \delta(s) + c$ by assumption.
We have $\dim_x(X_s) = \dim_{x'}(X_{s'})$ by
Morphisms, Lemma \ref{morphisms-lemma-dimension-fibre-after-base-change}.
Since $x$ is a generic point of an irreducible component $X_s$
(this is an irreducible scheme but we don't need this) and
$x'$ is a generic point of an irreducible component of $X'_{s'}$ we conclude
that $\text{trdeg}_{\kappa(s)}(\kappa(x)) = 
\text{trdeg}_{\kappa(s')}(\kappa(x'))$
by Morphisms, Lemma \ref{morphisms-lemma-dimension-fibre-at-a-point}.
Then
$$
\delta_{X'/S'}(x') = \delta(s') + \text{trdeg}_{\kappa(s')}(\kappa(x')) =
\delta(s) + c + \text{trdeg}_{\kappa(s)}(\kappa(x)) = \delta_{X/S}(x) + c
$$
This proves what we want by Definition \ref{definition-delta-dimension}.
\end{proof}

\end{document}
